% Results 4.3-4.5: top tagging, quark/gluon discrimination (Pythia and Herwig), fine-tuning with few
% training jets. Outline: docs/PAPER_OUTLINE.md sections 4.3-4.5; figures F5, F7; tables T3, T4.
% Every number is a macro from experiments/FIGS/make_tables.py (scratchpad phase_b/MACROS.md).

\subsection{Top tagging}\label{sec:top}

Top tagging uses the landscape dataset of hadronic top-quark jets and QCD jets~\cite{toptag} with its
standard training, validation and test sets, whose sizes are given with table~\ref{tab:benchtop}. The
dataset provides only the particle four-momenta, from which the kinematic inputs are computed. The
charge, the five particle-type flags and the four displacement variables, \FillTopNZero{} of the
\FillNFeatures{} inputs of the network, are therefore set to zero for every model. Published results are quoted
as their authors report them. The ROC curve of every label set after fine-tuning, on top tagging and on
quark/gluon discrimination, is drawn in figure~\ref{fig:f6}, and $R_{30}$ against the number of classes in
figure~\ref{fig:f2b}(c) and~(d).

\providecommand{\TableContinued}[2]{\label{#2}\ifnum\getpagerefnumber{#2}>\getpagerefnumber{#1}\relax{\footnotesize\textit{Table~\ref{#1} (continued)}}\par\smallskip\fi}
\begin{center}
\small
\begin{minipage}{\linewidth}\centering
\captionof{table}{Top tagging: Sophon's network pretrained on JetClass-II with 188, 162, 43 and 17 classes, or not pretrained, fine-tuned on the full top training set (\LitPartTopNTrainMillion{} million jets): accuracy, AUC and the background rejections $R_{50}$ and $R_{30}$, $1/\epsB$ at 50\% and 30\% signal efficiency, and the number of trainable parameters; (f.t.): fine-tuned. Published rows from the tables cited, the parameters from each model's paper (ParT and ParticleNet-f.t.: ParT's JetClass models; linear EFP: a linear classifier on energy flow polynomials); f.t.\ marks a model pretrained on JetClass and fine-tuned. Ours: \MlstNParamsFt{} trainable parameters, \MlstNParamsFtKin{} with kinematic inputs only. Our pretrained models see the \FillTopNZero{} inputs the top data lack set to zero; the model with kinematic inputs only, a diagnostic of that mismatch, has none missing (one run). Uncertainties: ours, the standard deviation over the pretraining runs, one fine-tuning each, and over three seeds from scratch; ParT, the spread over \LitPartFtHeadReinits{} fine-tunings of one pretrained model with different output-layer initializations, quoted for the run of median accuracy; L-GATr and OmniLearned, the standard deviation of five trainings; the other published uncertainties as given there. Bold: the best value of each column, and every value within one standard deviation of it (the larger of the two values' standard deviations; a configuration with two runs is compared by their mean and standard deviation).}
\label{tab:benchtop}
\par\medskip
\begin{adjustbox}{max width=\linewidth, max totalheight=0.55\textheight}
\begin{tabular}{>{\raggedright\arraybackslash}p{5.9cm}cc}
\toprule
 & \# params & accuracy \\
\midrule
linear EFP~\cite{toptag} & 1\,k & 0.932 \\
ParticleNet~\cite{particlenet} & 366\,k & 0.940 \\
ParT~\cite{part} & 2.14\,M & 0.940 \\
ParT-f.t.~\cite{part} & 2.14\,M & 0.944 \\
MIParT-L-f.t.~\cite{mipart} & 2.38\,M & 0.944 \\
L-GATr-f.t.~\cite{lgatr} & 1.1\,M & \textbf{0.9446} \\
OmniLearned-l~\cite{omnilearned} & 423\,M & 0.944 \\
\midrule
\multicolumn{3}{l}{\textit{this work}} \\
\mbox{no pretraining (f.t.)} & 2.21\,M & 0.9387\,$\pm$\,0.0002 \\
\mbox{\Kfull{} classes (f.t.)} & 2.21\,M & 0.9400\,$\pm$\,0.0002 \\
\mbox{\Kres{} classes (f.t.)} & 2.21\,M & 0.9401\,$\pm$\,0.0002 \\
\mbox{\Kmid{} classes (f.t.)} & 2.21\,M & 0.9402\,$\pm$\,0.0004 \\
\mbox{\Kcoarse{} classes (f.t.)} & 2.21\,M & 0.9398\,$\pm$\,0.0001 \\
\mbox{\Kfull{} classes,} \mbox{kinematic inputs only (f.t.)} & 2.21\,M & 0.9407 \\
\bottomrule
\end{tabular}
\end{adjustbox}
\end{minipage}
\par\bigskip
\begin{minipage}{\linewidth}\centering
\TableContinued{tab:benchtop}{tab:benchtop:2}
\begin{adjustbox}{max width=\linewidth, max totalheight=0.92\textheight}
\begin{tabular}{>{\raggedright\arraybackslash}p{5.9cm}ccc}
\toprule
 & AUC & $R_{50}$ & $R_{30}$ \\
\midrule
linear EFP~\cite{toptag} & 0.980 & -- & 384 \\
ParticleNet~\cite{particlenet} & 0.9858 & 397\,$\pm$\,7 & 1{,}615\,$\pm$\,93 \\
ParT~\cite{part} & 0.9858 & 413\,$\pm$\,16 & 1{,}602\,$\pm$\,81 \\
ParT-f.t.~\cite{part} & 0.9877 & \textbf{691\,$\pm$\,15} & 2{,}766\,$\pm$\,130 \\
MIParT-L-f.t.~\cite{mipart} & 0.9878 & 640\,$\pm$\,10 & 2{,}789\,$\pm$\,133 \\
L-GATr-f.t.~\cite{lgatr} & 0.9879 & 651\,$\pm$\,11 & 2{,}894\,$\pm$\,84 \\
OmniLearned-l~\cite{omnilearned} & \textbf{0.9880} & \textbf{688\,$\pm$\,9} & \textbf{3{,}486\,$\pm$\,157} \\
\midrule
\multicolumn{4}{l}{\textit{this work}} \\
\mbox{no pretraining (f.t.)} & 0.9852\,$\pm$\,0.0001 & 358\,$\pm$\,11 & 1{,}410\,$\pm$\,20 \\
\mbox{\Kfull{} classes (f.t.)} & 0.9861\,$\pm$\,0.0001 & 461\,$\pm$\,4 & 1{,}800\,$\pm$\,60 \\
\mbox{\Kres{} classes (f.t.)} & 0.9861\,$\pm$\,0.0001 & 443\,$\pm$\,17 & 1{,}790\,$\pm$\,60 \\
\mbox{\Kmid{} classes (f.t.)} & 0.9861\,$\pm$\,0.0001 & 450\,$\pm$\,13 & 1{,}780\,$\pm$\,150 \\
\mbox{\Kcoarse{} classes (f.t.)} & 0.9859\,$\pm$\,0.0001 & 429\,$\pm$\,10 & 1{,}740\,$\pm$\,140 \\
\mbox{\Kfull{} classes,} \mbox{kinematic inputs only (f.t.)} & 0.9862 & 443 & 1{,}800 \\
\bottomrule
\end{tabular}
\end{adjustbox}
\end{minipage}
\par\smallskip{\footnotesize ParT's top-tagging training script, unchanged, on our cluster gives $R_{30}$ of \MlstPartreproRthirtyPart{} (published 1{,}602\,$\pm$\,81); with Sophon's output head, \MlstPartreproRthirtyHead{}; with the head and Sophon's jet-$p_{\mathrm{T}}$ rescaling, \MlstPartreproRthirtyRescale{} (one run each; our network from scratch, which uses the rescaling, \MlstFtTopRthirtyScratch{}); table~\ref{tab:partrepro}.}
\end{center}


\begin{figure}[t]
\centering
\paperfig{\linewidth}{F6}
\caption{Fine-tuned models on the public benchmarks: background rejection $1/\epsB$ against signal efficiency
$\epsS$ on (a) top tagging, (b) quark/gluon discrimination on Pythia jets, with quark jets as the signal, and
(c) the same networks scored on Herwig jets, each with the two-class output of the network as the
discriminant, equation~(\ref{eq:discr}). Lines: the \Kfull-, \Kres-, \Kmid- and \Kcoarse-class models,
fine-tuned on the full training sets with the missing inputs set to zero, and the same network trained from
scratch, each the mean of its runs at fixed $\epsS$ (\MlstNRunsWord{} pretraining runs, one fine-tuning each;
\MlstFtRefSeedsWord{} training seeds from scratch); black dashes, in (a) only: the \Kfull-class model with
kinematic inputs only, pretrained once. Below each panel, each curve divided by the \Kfull-class curve. Each
curve stops where fewer than $\nmin$ background jets pass, equation~(\ref{eq:cut}). The values at the working
points are in tables~\ref{tab:benchtop} and~\ref{tab:benchqg}.}
\label{fig:f6}
\end{figure}

Pretraining raises the background rejection on top tagging, and no label set changes it beyond what
the test set resolves. Fine-tuned with the missing inputs set to zero, the \Kfull-class model reaches
\MlstFtTopRthirtyFullOverScratchPr{} than the same network trained from scratch ($R_{30}$), and every % CHECK: MlstFtTopRthirtyFullOverScratch above 1 by more than 2 SD (today 1.27)
label set improves on training from scratch in every run (table~\ref{tab:benchtop}). The % CHECK: MlstFtTopRthirty{Res,Mid,Coarse} every run above every from-scratch seed (today min 1,590 vs max 1,432)
\Kcoarse-class model matches the \Kfull-class model (\MlstFtTopRthirtyCoarseOverFullPr{}), as do the % CHECK: MlstFtTopRthirtyCoarseOverFullPr must read "equal within" (today 0.97 +- 0.10)
\Kres- and \Kmid-class models, so unlike the \Kcoarse-class labels on JetClass % CHECK: MlstFtTopRthirty{Res,Mid}OverFullPr must read "equal within"
(section~\ref{sec:jetclass}), no label set differs in $R_{30}$ on top tagging by the criterion of
section~\ref{sec:uncertainties}. Away from that working point the \Kcoarse-class model lies a few per cent
below the \Kfull-class model over much of the curve (figure~\ref{fig:f6}(a);
\MlstFtTopRfiftyCoarseOverFullPr{} at 50\% signal efficiency, comparable to the counting uncertainty there).
% CHECK: in F6(a) the ratio bands of 162/43/17 include one at eps_S = 0.3 (the R30 criterion).
% CHECK: MlstFtTopRfiftyCoarseOverFull (today 0.932 +- 0.018) near sqrt(1/n_17 + 1/n_188) (about 438 and 470 passing jets: 6.6%).
The test set may limit this comparison: $R_{30}$ is set by the small
number of QCD jets that pass at that working point, and their counting uncertainty, taken as independent for
the two models, is larger than the spread between runs (table~\ref{tab:boot}), although the shared test jets
make that figure an overestimate (appendix~\ref{sec:app-boot}). The agreement has a physical reason. A top-quark jet
differs from a QCD jet by its three-prong substructure and its mass, which every label set rewards
because even the \Kcoarse-class labels separate decays by their number of prongs. The pairwise
features of the network are computed from the four-momenta alone and do not depend on the missing
inputs.

Our network trained from scratch rejects less background than ParT trained from
scratch~\cite{part}, and the difference comes from Sophon's preprocessing, which every pretrained
model keeps. Run unchanged on our hardware, ParT's training script reproduces ParT's published
accuracy and AUC, and replacing ParT's output head with Sophon's leaves the rejection within the
published spread. Adding Sophon's rescaling of every jet's constituents to one fixed jet transverse
momentum~\cite{sophoncode} gives \MlstPartreproRthirtyRescaleOverPartPr{} than the unchanged script
(one run each; table~\ref{tab:partrepro}). We keep the rescaling in the network trained from scratch
because the pretrained models use it, so the comparison of pretraining with training from scratch does
not also change the preprocessing. The rescaling keeps the ratio of jet mass to transverse momentum but no longer the
jet mass itself, so the spread of transverse momenta in the sample likely broadens the top-quark and
W-boson mass peaks a top tagger relies on.

ParT fine-tuned after pretraining on JetClass~\cite{part} rejects more background than any of our
models, and the \Kfull-class model with kinematic inputs only measures how much of that gap the
missing inputs cause. ParT pretrained a separate model on the kinematic features for top tagging, so
nothing was missing for it. Our \Kfull-class model, with \FillTopNZero{} inputs set to zero, reaches % CHECK: MlstFtTopRthirtyFullOverPart below 1 (today 0.65)
\MlstFtTopRthirtyFullOverPartPr{} than ParT. The \Kfull-class model with kinematic inputs only
shares the data, the rescaling and the training settings of the \Kfull-class model and lacks no
input. Fine-tuned on top tagging, it leaves the gap to ParT as it is (kinematic-only model over the % CHECK: MlstFtTopRthirtyKinOverFull within the spread of the three 188-class runs
\Kfull-class model: \MlstFtTopRthirtyKinOverFullPr{}). It is one pretraining run, so it is read
against the spread of the \MlstNRunsWord{} \Kfull-class runs rather than as a ratio paired by seed. After fine-tuning the gap therefore lies in what else separates
our models from ParT: the rescaling above, and pretraining on JetClass-II, with pileup, rather than
on JetClass.

On top tagging, by the rule of section~\ref{sec:uncertainties}, frozen features of the four levels of the class
hierarchy give a linear probe no more than the randomly initialized network gives it. With the missing inputs set to zero, a linear probe on the frozen
\Kfull-class features and one on the randomly initialized network reach a similar rejection % CHECK: MlstLpTopRthirtyFullOverInitBnZeroPr reads "equal within" (two init_bn seeds: the Methods two-run rule; today 1.04 and 1.08); a stated gain contradicts "no more"
(the zero column of top tagging in the lower part of table~\ref{tab:fill}; \Kfull-class over randomly initialized network: \MlstLpTopRthirtyFullOverInitBnZeroPr{}; the MLP probe:
\MlstMlpTopRthirtyFullOverInitBnPr{}), and both lie far below every fine-tuned model
(figures~\ref{fig:f2a}(c) and~\ref{fig:f2b}(c) and table~\ref{tab:benchtop}). The auxiliary mass regression can add to this: the \Kres-class % CHECK: MlstLpTopRthirty{Res,Mid,Coarse}OverInitBnZeroPr read "equal within" (today 40, 30, 20%) and MlstLpTopRthirtyResMassOverInitBnZeroPr states a gain (today at least 68%); MlstLpTopRthirtyCoarseMassOverInitBnZero has one run near 1 and one far above (today 2.3 and 1.2); recheck with the substitution chosen on validation (fits_v4_fill)
model with it gains in both runs (\MlstLpTopRthirtyResMassOverInitBnZeroPr{}), and the two runs of the \Kcoarse-class
model with it differ too much to read (\MlstLpTopRthirtyCoarseMassOverInitBnZero{} times the randomly initialized network). Fine-tuning the % CHECK: MlstMlpTopRthirtyFullOverInitBnPr must read "equal within" (today 1.00 +- 0.31)
same backbone on the same zero-substituted jets gives the gain of table~\ref{tab:benchtop}, so the
frozen features are limited by how they respond to the substituted inputs, not by what the backbone
can represent.

\begin{landscape}
\begin{table*}[t]
\centering
\small
\caption{Linear probes on the frozen features of networks pretrained on JetClass-II with 188, 162, 43 and 17 classes, with the inputs a data set lacks set to zero, to the training mean or to random draws of a real particle, trained at the largest training set; (LP): linear probe on the frozen features. Upper part: relative drop in JetClass-II accuracy (percent) when the slots the top or the quark/gluon data lack are substituted, against the same model with every input. Lower part: $R_{30}$ on top and quark/gluon tagging. Mean and SD over three runs. The \Kfull-class model with kinematic inputs only has nothing missing (one run). The controls are in table~\ref{tab:fill}.}
\label{tab:fill_main}
\begin{tabular*}{\linewidth}{@{\extracolsep{\fill}}lcccccc@{}}
\toprule
 & \multicolumn{3}{c}{JetClass-II, top slots} & \multicolumn{3}{c}{JetClass-II, quark/gluon slots} \\
\cmidrule(lr){2-4}\cmidrule(lr){5-7}
 & zero & mean & draws & zero & mean & draws \\
\midrule
\Kfull{} classes (LP) & 88.9\,$\pm$\,1.0\% & 87.9\,$\pm$\,1.1\% & 85.98\,$\pm$\,0.07\% & 61.2\,$\pm$\,3.0\% & 48.8\,$\pm$\,1.1\% & 54.3\,$\pm$\,0.2\% \\
\Kres{} classes (LP) & 89.0\,$\pm$\,0.9\% & 87.8\,$\pm$\,0.7\% & 85.89\,$\pm$\,0.11\% & 64.1\,$\pm$\,5.2\% & 49.5\,$\pm$\,0.5\% & 54.6\,$\pm$\,0.6\% \\
\Kmid{} classes (LP) & 87.8\,$\pm$\,0.3\% & 86.7\,$\pm$\,0.9\% & 83.17\,$\pm$\,0.29\% & 67.6\,$\pm$\,4.8\% & 47.6\,$\pm$\,1.6\% & 53.5\,$\pm$\,1.4\% \\
\Kcoarse{} classes (LP) & 86.7\,$\pm$\,0.7\% & 86.1\,$\pm$\,0.8\% & 81.42\,$\pm$\,0.53\% & 70.9\,$\pm$\,2.3\% & 50.0\,$\pm$\,1.0\% & 51.7\,$\pm$\,0.4\% \\
\bottomrule
\end{tabular*}
\par\bigskip
\begin{tabular*}{\linewidth}{@{\extracolsep{\fill}}lcccccc@{}}
\toprule
 & \multicolumn{3}{c}{top, $R_{30}$} & \multicolumn{3}{c}{quark/gluon, $R_{30}$} \\
\cmidrule(lr){2-4}\cmidrule(lr){5-7}
 & zero & mean & draws & zero & mean & draws \\
\midrule
\Kfull{} classes (LP) & 60\,$\pm$\,12 & 73\,$\pm$\,15 & 228\,$\pm$\,25 & 66\,$\pm$\,7 & 72.7\,$\pm$\,2.4 & 68.1\,$\pm$\,1.8 \\
\Kres{} classes (LP) & 76\,$\pm$\,17 & 64\,$\pm$\,11 & 199\,$\pm$\,25 & 58\,$\pm$\,9 & 71.0\,$\pm$\,0.3 & 66.8\,$\pm$\,0.4 \\
\Kmid{} classes (LP) & 57\,$\pm$\,13 & 58\,$\pm$\,14 & 231\,$\pm$\,8 & 61\,$\pm$\,9 & 70.8\,$\pm$\,1.8 & 69.0\,$\pm$\,1.1 \\
\Kcoarse{} classes (LP) & 55\,$\pm$\,5 & 55\,$\pm$\,5 & 255\,$\pm$\,17 & 62\,$\pm$\,7 & 74.3\,$\pm$\,1.9 & 71.6\,$\pm$\,2.8 \\
\Kfull{} classes, kinematic inputs only (LP) & \multicolumn{3}{c}{231 (nothing missing)} & -- & -- & -- \\
\bottomrule
\end{tabular*}
\end{table*}
\end{landscape}


The same zeros cost the pretrained features on JetClass-II itself, more than they cost the randomly
initialized network. With the inputs the top data lack set to zero in the training and test jets of
the linear probe, the \Kfull-class model loses \MlstLpJcTwoTopZeroAccDropFullPr{} of its accuracy on
the \Kfull{} classes and the randomly initialized network \MlstLpJcTwoTopZeroAccDropInitBnPr{}. Every % CHECK: MlstLpJcTwoTopZeroAccDrop{Res,Mid,Coarse,ResMass,CoarseMass,Fzero,Fone,Foner} each above MlstLpJcTwoTopZeroAccDropInitBn (today 85-90 vs 55-58)
pretrained label set loses more than the randomly initialized network (the top-slot columns of
tables~\ref{tab:fill_main} and~\ref{tab:fill}). Two explanations predict this pattern. The zeroed inputs are states the
pretrained network never saw, so its features respond to them in ways no training jet shaped.
Alternatively, the type flags and displacements carry information that the pretrained features use and
the random features do not, so the pretrained network has more to lose. The other two substitutions separate
these explanations.

The substitutions differ in how far they leave each particle from a real one, and that difference
separates the two explanations. No JetClass-II particle has all of its type flags at zero, so the
zeroed inputs put every particle of a top jet in a joint state the pretrained networks never saw. Each zeroed flag on its own still lies
within the range of the training distribution after the input normalization (table~\ref{tab:inputs});
it is the combination that no real particle has. The training mean is no closer to a real particle, since it is the centroid
of the training particles and carries fractional flags; only each random draw is a state a real
particle takes. The top-tagging columns in the lower part of table~\ref{tab:fill_main} compare the linear
probe on the \Kfull-class model with random draws and with zeros (draws over zeros: \MlstLpTopRthirtyDrawOverZeroFullPr{}), and the linear
probe on the \Kfull-class model with kinematic inputs only, which lacks nothing, with the
zero-substituted \Kfull-class model (\MlstLpTopRthirtyKinOverFullPr{}). The draws bring the probe % CHECK: draws close to kinematic-only, one run (both ratios 3.9 at c0ce260); "close to", not "the same": the two agree on R30 and less exactly on AUC and R50
close to the model that lacks nothing, which places the loss in the zero value rather than in the missing information.
Fine-tuning escapes both, because it re-learns the input layer on
the substituted jets, and a constant substitution used in training and testing is consistent for a
model trained with it~\cite[Theorem~3]{josse}. JetParticle-JEPA, which saw zeroed inputs during pretraining through
input dropout, sets the missing top-tagging inputs to zero in the same way and reports a rejection
below that of ParT trained from scratch~\cite{jpjepa}; its result does not distinguish the two
explanations either.

\subsection{Quark/gluon discrimination, Pythia and Herwig}\label{sec:qg}

Quark/gluon discrimination uses the Pythia sample of Komiske et al.~\cite{efn,qgpythia} with ParT's
split into training, validation and test jets~\cite{part}, and the companion sample generated with
Herwig with the same selection~\cite{herwigsevenone,qgherwig}; the sizes and cuts are given with
table~\ref{tab:benchqg}. The jets carry the particle type as five binary flags and the charge, ParT's
``exp'' inputs~\cite{part}, and only the \FillQgNZero{} displacement inputs are set to zero. Every
model is fine-tuned on Pythia jets and tested on both samples. Published results are from the same
input scenario.

\begin{table*}[t]
\centering
\small
\setlength{\tabcolsep}{3pt}
\caption{Quark/gluon tagging: Sophon's network pretrained on JetClass-II with 188, 162, 43 and 17 classes and fine-tuned, or trained from scratch, on the full training set (\LitQgTrainMillion{} million Pythia jets; particle type as five flags and the charge) and scored on Pythia jets, beside published models: accuracy, AUC, $1/\epsB$ at 50\% and 30\% quark efficiency and the number of trainable parameters. Uncertainties as in table~\ref{tab:benchtop}; Herwig jets in table~\ref{tab:benchqgherwig}.}
\label{tab:benchqg}
\begin{tabular*}{\linewidth}{@{\extracolsep{\fill}}>{\raggedright\arraybackslash}p{2.9cm}ccccc@{}}
\toprule
 & \# params & accuracy & AUC & $R_{50}$ & $R_{30}$ \\
\midrule
constituent multiplicity~\cite{efn} & -- & -- & 0.8401 & 19.0 & -- \\
linear EFP~\cite{efn} & -- & -- & 0.8919 & 29.7 & -- \\
PFN-Ex~\cite{efn} & -- & -- & 0.9005 & 34.7\,$\pm$\,0.4 & -- \\
ParticleNet~\cite{particlenet} & 366\,k & 0.840 & 0.9116 & 39.8\,$\pm$\,0.2 & 98.6\,$\pm$\,1.3 \\
ParT~\cite{part} & 2.14\,M & 0.840 & 0.9121 & 41.3\,$\pm$\,0.3 & 101.2\,$\pm$\,1.1 \\
ParticleNet-f.t.~\cite{part} & 370\,k & 0.839 & 0.9115 & 40.1\,$\pm$\,0.2 & 100.3\,$\pm$\,1.0 \\
ParT-f.t.~\cite{part} & 2.14\,M & 0.843 & 0.9151 & 42.4\,$\pm$\,0.2 & 107.9\,$\pm$\,0.5 \\
OmniLearned-m~\cite{omnilearned} & 58\,M & 0.845 & 0.9162 & 43.2\,$\pm$\,0.1 & 111.2\,$\pm$\,1.5 \\
\midrule
\multicolumn{6}{l}{\textit{this work}} \\
from scratch & 2.21\,M & 0.8397\,$\pm$\,0.0007 & 0.9121\,$\pm$\,0.0002 & 41.47\,$\pm$\,0.29 & 104.0\,$\pm$\,1.0 \\
\multicolumn{6}{l}{\textit{fine-tuned}} \\
\Kfull{} classes & 2.21\,M & 0.8433\,$\pm$\,0.0002 & 0.9148\,$\pm$\,0.0001 & 42.56\,$\pm$\,0.10 & 107.3\,$\pm$\,2.2 \\
\Kres{} classes & 2.21\,M & 0.8432\,$\pm$\,$<$0.0001 & 0.9148\,$\pm$\,0.0001 & 42.39\,$\pm$\,0.22 & 108.0\,$\pm$\,1.4 \\
\Kmid{} classes & 2.21\,M & 0.8430\,$\pm$\,0.0005 & 0.9146\,$\pm$\,0.0003 & 42.25\,$\pm$\,0.77 & 108.3\,$\pm$\,1.0 \\
\Kcoarse{} classes & 2.21\,M & 0.8427\,$\pm$\,0.0003 & 0.9144\,$\pm$\,$<$0.0001 & 41.97\,$\pm$\,0.25 & 106.5\,$\pm$\,0.8 \\
\bottomrule
\end{tabular*}
\end{table*}


Fine-tuned on the Pythia sample, the pretrained models reach the rejection of ParT fine-tuned after
pretraining on JetClass~\cite{part}, and the label sets do not differ. The \Kfull-class model
matches ParT (\MlstFtQgRthirtyFullOverPartPr{}), and the \Kcoarse-class model matches the % CHECK: MlstFtQgRthirtyFullOverPartPr must read "equal within" (today 0.99 +- 0.02)
\Kfull-class model (\MlstFtQgRthirtyCoarseOverFullPr{}; table~\ref{tab:benchqg}; figure~\ref{fig:f6}(b)), although only the % CHECK: MlstFtQgRthirtyCoarseOverFullPr must read "equal within" (today 0.99 +- 0.03); also Res, Mid
\Kfull- and \Kres-class labels separate resonance decays to gluons from decays to light quarks, the
distinction quark/gluon discrimination asks for. Three differences between the tasks likely limit how
much of that distinction transfers: the quark/gluon jets have a smaller radius than every pretraining jet; they
are isolated jets from a single parton rather than prongs inside a resonance decay; and their
displacement inputs are zero.

The change of generator costs the pretrained and the from-scratch networks the same fraction of their
rejection, so pretraining gives no protection against it. Tested on the Herwig sample, the fine-tuned
\Kfull-class model keeps \MlstFtQgRthirtyHerwigOverPythiaFullPr{} it has on Pythia jets, and the % CHECK: MlstFtQgRthirtyHerwigOverPythiaFull vs ...Scratch must agree within 2 SD (today 0.33 +- 0.02 vs 0.34 +- 0.01)
network trained from scratch \MlstFtQgRthirtyHerwigOverPythiaScratchPr{} (table~\ref{tab:benchqg};
figure~\ref{fig:f6}(c)).
Earlier generator comparisons found Herwig's quark and gluon jets less separable than
Pythia's~\cite{gras}, which lowers the rejection of any tagger trained on Pythia jets. Pretraining
does not offset this because JetClass-II is also showered and hadronized with
Pythia~\cite{sophon,pythia}, so the pretrained network has seen no jet with Herwig's modeling.

On quark/gluon discrimination, frozen features from pretraining give a linear probe more than the
randomly initialized network gives it, but not much more. With the displacement inputs substituted as
chosen on validation (\LpSubstitutionQg{}), a linear probe on the frozen \Kfull-class features
rejects more background than one on the randomly initialized network
(\MlstLpQgRthirtyFullOverInitBnPr{}; the quark/gluon columns in the lower part of table~\ref{tab:fill}). The frozen features likely do better
here than on top tagging because the type flags and the charge are present. They gain little over the
randomly initialized network because the missing displacements still leave every charged particle in
a state no pretraining jet has, as the next paragraph explains.

Setting only the displacement inputs to zero already costs the pretrained features most of their % CHECK: MlstLpJcTwoQgZeroAccDropFull above 50 ("most"; today 61)
accuracy on JetClass-II, and the randomly initialized network little. With the \FillQgNZero{} inputs the
quark/gluon data lack set to zero in the linear probe's training and test jets, the \Kfull-class
model loses \MlstLpJcTwoQgZeroAccDropFullPr{} of its accuracy on the \Kfull{} classes and the
randomly initialized network \MlstLpJcTwoQgZeroAccDropInitBnPr{} (the quark/gluon-slot columns of
table~\ref{tab:fill}). The % CHECK: MlstLpJcTwoQgZeroAccDropFull well above ...InitBn (today 61 vs 9.7-11)
mismatch differs from top tagging in one respect: a zeroed impact parameter lies at the training mean,
but no charged JetClass-II track has zero measurement uncertainty on its impact parameters, so a
charged particle with zeroed uncertainties is a joint state the input layer never saw. As on top
tagging, the pretrained network may also lose information it uses, since the displacements mark the
b- and c-hadron decays of many JetClass-II classes~\cite{cmsbtag}.

\subsection{Fine-tuning with few training jets}\label{sec:fewjets}

\pending{How the gain of the fine-tuned \Kfull-class model over training from scratch on top tagging
changes with the number of training jets, the data-efficiency equivalence of section~\ref{sec:evaluation},
and the order of the label sets at the smaller training sets (table~\ref{tab:ft}) are written once the
fine-tuning runs on the smaller training sets are evaluated.}
% CHECK: pending MlstFtTopRthirtyFullOverScratchEThree, MlstFtTopDataEquivJets (PI 2026-11-02: nothing is written about the top-tagging gain against training-set size until the fixed-budget reruns read out; the writer's draft claimed "pretraining helps most with few jets", which needs EThree above the full-set MlstFtTopRthirtyFullOverScratch)
% CHECK: pending -- the label-set order at 1e3-1e5 jets (table tab:ft, figure fig:a10) must match the full set; otherwise state the small-size difference with its ratio

\begin{figure}
\centering
\paperfig{\linewidth}{F8}
\caption{Pretraining against the number of training jets: (a) accuracy on the \Kfull{} classes of
class-balanced JetClass-II and (b) on JetClass, and $R_{30}$, the background rejection at 30\% signal
efficiency, on (c) top tagging and (d) quark/gluon discrimination with the Pythia sample, for the \Kfull-class
model fine-tuned and with a linear probe on its frozen features, and for the same network trained from
scratch. Points are means and bands the standard deviation over runs. Stars are ParT's published results: with
the same inputs on JetClass, kinematic inputs only on top tagging, and kinematic inputs with the particle type
on quark/gluon discrimination. On top and quark/gluon discrimination the linear probe's missing inputs are
substituted as chosen on validation (\LpSubstitutionTop{} and \LpSubstitutionQg{}), and fine-tuning sets them
to zero. Every label set at every size is in table~\ref{tab:ft} and figure~\ref{fig:a10}.}
\label{fig:f8}
\end{figure}

Most of the gain from pretraining survives with the backbone frozen when no input is missing. On the
smallest JetClass training set the linear probe on the \Kfull-class model keeps
\MlstLpJcOneGapFracEFourPr{} of the gain of fine-tuning over training from scratch
(section~\ref{sec:evaluation}), and on class-balanced
JetClass-II \MlstLpJcTwoGapFracEFourPr{}. On top tagging and quark/gluon discrimination with the full
training sets, the linear probe stays far below fine-tuning (figure~\ref{fig:f8}(c, d); tables~\ref{tab:probesn}
and~\ref{tab:ft}), because fine-tuning re-learns the input layer on the zeroed inputs, while the frozen
features keep the mismatch described in section~\ref{sec:top}.
% CHECK: pending the final budget read-out -- at 1e3 top jets the 188-class probe (random draws, R30 162+-9) exceeds
% fine-tuning (zero fill, 51/78/87) and from scratch (45-55); from 1e4 fine-tuning is above (about 420 vs the probe). Write the
% size dependence in sec:fewjets (PI 2026-11-02 rule), naming that probe and fine-tuning fill the missing inputs differently.
